\subsection*{从直角三角形走到整圈运动}
锐角三角比只覆盖部分角度，单位圆则能连续描述整圈。点在圆上转到角度 $\theta$ 时，横坐标定义为余弦，纵坐标定义为正弦。负角、超过一圈的角和各象限符号，都由点的位置统一解释。

半径为一，所以两坐标平方和为一。这不是一条需要孤立记住的恒等式，而是点始终在单位圆上的坐标表达。关于纵轴或横轴反射，也会直接给出一批符号关系。

\subsection*{和角公式来自两次旋转}
先旋转 $\beta$，再旋转 $\alpha$，总角度是 $\alpha+\beta$。把每次旋转写成矩阵，乘积的第一列记录标准横向箭头最终落点；与总旋转矩阵比较，就得到正弦余弦的和角公式。

这条路线解释了为什么正弦和角不是简单相加：第二次旋转会同时混合前一步的横纵分量。公式中的乘积项记录这种投影与混合，而非人为增加的复杂性。

在复平面实验里，先暂停在某个角度，预测再转四分之一圈后的横纵坐标，再改变滑块验证。学习三角时应能在角度、点坐标、旋转动作三种语言之间转换。后面复数乘法会把同一个旋转矩阵压缩成一个复数。
